The geometry of Section~\ref{sec:p0-aes} is developed on the basic model, which
is complete and unpunctured: $\mathfrak E_0$ has no point at which the
arithmetic structure fails.  (The one-puncture model $\mathfrak E_1$ of
Example~\ref{ex:p0-e1-model} already leaves the regular category, and is the
case this chapter generalizes.)  This final chapter records the opposite
situation, which arose in the most recent
exploratory work of the programme, and the vocabulary that has been proposed for
it.  The mathematical content of the chapter is deliberately modest: one
definition, one exact local computation, and a register of six computed
witnesses, each with its own status.  The interpretive reading of those
witnesses --- the language of tearing --- is presented as a structural proposal,
and Subsection~\ref{subsec:p0-holed-nonclaims} lists what is not established.

\subsection{Punctured arithmetic expression spaces}
\label{subsec:p0-holed-definition}

\begin{definition}[Punctured arithmetic expression space]
\label{def:p0-punctured-aes}
Let $\overline M$ be an oriented smooth surface, let $S\subset\overline M$ be a
finite set of points, and put $M=\overline M\setminus S$.  Suppose given a
smooth Riemannian metric $g$ on $M$, a function $a:M\to\mathbb R$ that is smooth
on $M$, and constants $\mu\ne0$, $\lambda$, such that
\begin{enumerate}
  \item $(M,g,a;\mu,\lambda)$ is a regular AES in the sense of
  Definition~\ref{def:regular-aes}, and
  \item for every $p\in S$, no neighbourhood $U$ of $p$ in $\overline M$ admits
  extensions of $g|_{U\cap M}$ and $a|_{U\cap M}$ to smooth data
  $\widetilde g,\widetilde a$ on $U$ for which
  $(U,\widetilde g,\widetilde a;\mu,\lambda)$ is a regular AES.
\end{enumerate}
Then $(\overline M,S,g,a;\mu,\lambda)$ is called a \emph{punctured arithmetic
expression space}, and the points of $S$ are its \emph{punctures}.  When
$|S|=k$ we write $\mathfrak E_k$ for it, in the descriptive sense of
Convention~\ref{conv:p0-model-labels}.
\end{definition}

The ambient surface and the data are deliberately separated: the regular AES is
given only on the punctured surface $M$, so that condition (ii) is a genuine
requirement rather than a contradiction.  A puncture is not any point that has
been deleted, but a point at which the regular structure genuinely cannot be
continued.  This is the finite-point case of the singular arithmetic expression
space of Paper~I; the general singular theory, the classification of singular
germs, and everything concerning tubes, braids, and knot invariants belong to
Paper~I and Paper~III.  The notion is specialized here for one reason only: the
arithmetic motions of Subsection~\ref{subsec:p0-point-paths} are defined on the
punctured surface, and the question of what happens to them at a puncture is a
question about AES.

The first example is the one-puncture model of Example~\ref{ex:p0-e1-model}: the
unit disc with the complete Poincar\'e metric and the radial assignment
$a_{\mathbb D}=\frac{2\mu}{\lambda}\rho/(1-\rho^2)$, whose metric extends
smoothly across the centre while the assignment does not.  There $S=\{0\}$ and
the regular locus is $\mathbb D\setminus\{0\}$.  Two features of that example
are worth recording here, because they are typical of punctured models.  First,
the extended metric is complete on $\mathbb D$ while its restriction to
$\mathbb D\setminus\{0\}$ is \emph{not} complete: the puncture is at finite
distance from every interior point, so the incompleteness is caused by the
deletion and not by the geometry.  Second, the assignment has a finite limit at
the puncture (here $0$) but no $C^1$ extension, so the failure is a failure of
differentiability rather than of continuity.  Both statements are verified in
Paper~I together with the model.

\subsection{The local shape of a puncture}
\label{subsec:p0-holed-tangent-cone}

A puncture is not merely a missing point; at least in the example that has been
computed in detail, it carries a definite local shape.  The following statement
is exact.

\begin{proposition}[Tangent cone of the four limiting circles]
\label{prop:p0-puncture-tangent-cone}
Fix $\rho>0$ and consider the four circles through the origin
\[
  C_\pm:\ x^2+y^2\mp2\rho y=0,
  \qquad
  D_\pm:\ x^2+y^2\mp2\rho x=0 .
\]
Write $S_2=x^2+y^2$ and let
\[
  P=\bigl(S_2^2-4\rho^2y^2\bigr)\bigl(S_2^2-4\rho^2x^2\bigr)
\]
be the defining polynomial of their union.  Then
\[
  P=S_2^4-4\rho^2S_2^3+16\rho^4x^2y^2 ,
\]
so the lowest-degree homogeneous part of $P$ is $16\rho^4x^2y^2$.  Consequently
the tangent cone of the union at the origin is
\[
  (xy)^2=0 ,
\]
that is, the two coordinate axes, each with multiplicity two.
\end{proposition}

\begin{proof}
The four circles are $S_2=\pm2\rho y$ and $S_2=\pm2\rho x$, whose defining
polynomials are $S_2^2-4\rho^2y^2$ and $S_2^2-4\rho^2x^2$ respectively; the
product is $P$.  Expanding,
\[
  \bigl(S_2^2-4\rho^2y^2\bigr)\bigl(S_2^2-4\rho^2x^2\bigr)
  =S_2^4-4\rho^2S_2^2(x^2+y^2)+16\rho^4x^2y^2
  =S_2^4-4\rho^2S_2^3+16\rho^4x^2y^2 ,
\]
since $S_2^2(x^2+y^2)=S_2^3$.  The terms have degrees eight, six, and four, so
the lowest-degree part is $16\rho^4x^2y^2$, whose zero set is $xy=0$ counted
twice.
\end{proof}

Two remarks, both used below.  First, the two circles of the same family are
tangent at the origin, so the origin is a degenerate point of each family, while
circles of different families cross there with tangent lines $x=0$ and $y=0$;
the tangent cone is the union of those two lines.  Second, the tangent cone is
*reducible with multiplicity*: it is the square of a product of two distinct
linear forms.

No arithmetic expression space is attached to this configuration in this paper.
The four circles are an explicit curve configuration whose tangent cone can be
computed exactly, and nothing more is claimed: in particular they are not
asserted to be the limiting circles of any AES metric or assignment, and the
origin is not asserted to be a puncture in the sense of
Definition~\ref{def:p0-punctured-aes}.  Constructing data $(g,a)$ on a
neighbourhood of the origin for which this configuration is natural, and for
which the origin is non-extendable, is recorded as the first problem of
Subsection~\ref{subsec:p0-holed-open}.
The computation is recorded in the exploration register of the
programme \cite{Yuan2026AEGHoleRegister}; the identity above is verified here
directly.

\subsection{Tearing at a puncture}
\label{subsec:p0-tearing-at-puncture}

The tearing of Definition~\ref{def:p0-tearing} is a defect between two motions
of a space.  On a punctured space a further question arises: the motions
themselves need not be extendable across the puncture, so a comparison of two
motions may have to be performed around it rather than through it.  This
subsection records that question and what is known about it; it proves nothing
beyond the definitions.

Let $(\overline M,S,g,a;\mu,\lambda)$ be a punctured AES and let $p\in S$.  A
motion of the punctured surface is a path in $M=\overline M\setminus S$, and two
motions that approach $p$ from different directions need not have limits related
by any regular motion of the punctured part: the assignment and the metric are
defined only on $M$.  The following two questions are therefore well posed.

\begin{enumerate}
  \item \emph{Does comparison require a detour?}  Two bounded motions with the
  same charge whose paths cannot be deformed into one another in $M$ while
  keeping their charges: is their relative torsion still given by an ACS area,
  and if so in which filling class?  The frame witness of
  Subsection~\ref{subsec:p0-holed-witnesses} is the computed instance that makes
  the question concrete, because there a closed loop in one coordinate is open
  in another.
  \item \emph{Does comparison require a cut?}  If the punctures are stipulated
  rather than derived, then the family of comparisons depends on that
  stipulation.  The gluing witness of the same subsection exhibits a whole
  family of cuts whose invariants are exhausted by a permutation, so the cut
  carries more information than the topology it produces.
\end{enumerate}

Neither question is answered here.  What can be said now is only where the
measurement must be performed: the measured value of tearing is the ACS area of
Theorem~\ref{thm:torsion-stokes}, which is a statement about add--scale
histories and therefore does not itself refer to the puncture; a comparison
across a puncture has to leave the punctured surface, either by passing to the
charge plane or by going around the puncture, and deciding which of the two is
canonical is recorded as the second problem of
Subsection~\ref{subsec:p0-holed-open}.

\subsection{Six computed witnesses}
\label{subsec:p0-holed-witnesses}

The exploratory work recorded in \cite{Yuan2026AEGHoleRegister} produced six
separately computed statements about punctured configurations.  They are
reproduced here in the register's own order, each with its exact content and its
status.  None of them is promoted by appearing in this paper: the register
carries no authoritative status, and the items below are classified according to
the labels of this repository.

\begin{enumerate}
  \item \emph{A topological witness.}  Cut a disc carrying four right-angle
  sectors along the two diagonal seams and the four coordinate rays; eight
  triangular pieces result, each with a left radial edge, an outer arc, and a
  right radial edge.  Regluing is a permutation $\sigma$ of the eight left edges
  onto the eight right edges.  If $C$ is the number of cycles of $\sigma$, then
  the centre vertex has $C$ orbits, the outer vertices form eight orbits, and
  $V=C+8$, $E=16$, $F=8$, whence
  \[
    \chi=V-E+F=C=b ,
  \]
  the number of boundary components.  Consequently every gluing of this family
  produces either one disc or a disjoint union of $C$ discs, and the compact
  three-holed sphere ($\chi=-1$ with three boundary components) is unreachable
  in this family.  A surface of Euler characteristic $-1$ is obtained instead by
  removing two points from the disc: it has one boundary component, two
  punctures, and hence three ends, so its Euler characteristic agrees with that
  of the compact pair of pants while its boundary structure does not.  An
  exhaustive check of all $40320$ permutations confirms the
  distribution $5040,13068,13132,6769,1960,322,28,1$ for $C=1,\ldots,8$.
  \emph{Status:} the identity $\chi=C$ is elementary and provable as above; the
  enumeration is a computationally verified example.  The topological
  conclusion holds for the stated gluing rule, not for arbitrary gluings of the
  eight pieces.

  \item \emph{A combinatorial witness.}  If three components interact pairwise,
  the interaction sites are the three unordered pairs, and each site carries two
  directions of exchange, giving six directed transitions in total.
  \emph{Status:} counting, given the definition of a hole as an interaction
  site.  The definition is a choice, not a theorem.

  \item \emph{A frame witness.}  Let the state of a three-component frame be
  $(c,\sigma,\varphi)$ with component label $c\in\mathbb Z_3$, direction
  $\sigma\in\{\pm1\}$, and phase $\varphi\in\mathbb Z_4$, and let one step act by
  \[
    (c,\sigma,\varphi)\longmapsto(c+\sigma,\sigma,\varphi+\sigma).
  \]
  Since $\gcd(3,4)=1$, the Chinese remainder theorem identifies
  $(c,\varphi)$ with a single counter $k$ modulo $12$, and the orbit of each
  direction has length $12$.  Within one such orbit the component label
  advances in the direction $\sigma$, so the transitions realized are the three
  pairs of that direction, each occurring four times --- once per phase class.
  The six directed transitions of the combinatorial witness therefore appear
  four times each only over the two orbits together, that is, over $24$ steps;
  a single $12$-step orbit cannot contain $24$ transitions.  Since one full
  phase period is four steps and $4\equiv1\pmod3$, one period of the phase
  returns the phase to its initial value while advancing the component label by
  \[
    c\longmapsto c+\sigma \pmod 3 ;
  \]
  for the fixed direction $\sigma=+1$ this is the advance by one used in the
  register, and for $\sigma=-1$ it is the advance by two.  In either case the
  phase loop is not closed in the total state: its monodromy is the generator
  $-\sigma$ of $\mathbb Z_3$.  \emph{Status:} exact arithmetic (Chinese
  remainder theorem), verified exhaustively on the two $12$-state orbits and by
  replay.

  \item \emph{An observer witness.}  Let a frame be described by an angle
  $\theta$ and three joint offsets $\delta_1,\delta_2,\delta_3$, and let the
  residuals available to an observer be
  \[
    r_i=\bigl|2\sin\delta_i\bigr| \quad (i=1,2,3),
    \qquad
    r_4=\bigl|2\sin\bigl(2\theta+\tfrac12(\delta_1+\delta_2+\delta_3)\bigr)\bigr| .
  \]
  All four residuals vanish exactly at the ideal frame
  $(\theta,\delta)=(\pi/2,0,0,0)$.  Suppose an observer sees only residuals of
  magnitude at least $\varepsilon$, and that an iterative procedure stops as
  soon as every visible residual is below that threshold.  At the halt one then
  has the following bound.  From $r_i<\varepsilon$ we get
  $|\delta_i|<\arcsin(\varepsilon/2)=\varepsilon/2+O(\varepsilon^3)$, and from
  $r_4<\varepsilon$ together with
  $r_4=\bigl|4(\theta-\pi/2)+(\delta_1+\delta_2+\delta_3)\bigr|+O(\lVert(\theta-\pi/2,\delta)\rVert^2)$
  we get
  \[
    \bigl|\theta-\pi/2\bigr|
    <\frac58\,\varepsilon+O(\varepsilon^2).
  \]
  The halting offset is therefore of order $\varepsilon$ and decreases linearly
  in $\varepsilon$; it is not zero, and it is the best that a procedure obeying
  this stopping rule can report.  \emph{Status:} the residual functions, their
  exact common zero at the ideal frame, and the bound are proved above.  The
  particular iteration counts and offsets recorded in the register depend on the
  implementation of the update step, which the register does not fix, and they
  are accordingly \emph{not} reproduced here.  The witness is \emph{not} an
  obstruction: nothing in the computation shows that no procedure can do better,
  and the exact zero of the residuals at the ideal frame shows the opposite.

  \item \emph{A semantic witness.}  Model the three interaction sites by
  vocabularies built from four terms common to all three components, four
  further control terms distributed over the pairs as two, one, and one, and
  three terms private to the individual components (which therefore enter no
  site vocabulary).  The three site vocabularies then have sizes
  \[
    6=4+2,\qquad 5=4+1,\qquad 5=4+1 .
  \]
  Filling the three sites independently --- one word chosen from
  each site's own vocabulary --- gives $6\cdot5\cdot5=150$ possibilities;
  filling them so that the same core word is used at all three sites gives
  $4^3=64$.  The difference $86$ is carried by the adapters at the sites.  On
  the control layer the pairwise intersections are non-empty while the common
  intersection is empty; on the full vocabularies the common intersection is the
  four-term core, so the emptiness statement is layer-dependent.
  \emph{Status:} the arithmetic is verified given the model; the model is a
  modelling choice, and the "Borromean" description is an analogy with no
  link-theoretic content.

  \item \emph{A phase-geometry witness.}  With components at phases
  $\{0,+1,-1\}$ modulo $n$, a forward motion of $k$ steps and a reverse motion of
  $m$ steps meet exactly when $k+m\equiv0\pmod n$.  A meeting phase different
  from the three component phases exists precisely for $n\ge4$; at $n=4$ it is
  the first case in which the meeting is simultaneously balanced ($k=m$) and
  antipodal ($\varphi=n/2$), namely $k=m=2$, $\varphi=2$.  \emph{Status:}
  computationally verified example over the stated finite model, exhaustive for
  $n=3,\ldots,8$.  The value $n=4$ is minimal, not unique: $n=6$ and $n=8$ also
  admit balanced antipodal meetings.
\end{enumerate}

\begin{remark}[Register bookkeeping]
\label{rem:p0-holed-register-bookkeeping}
The exploration register introduces these items as five classes while tabulating
six rows; the sixth row is the phase-geometry witness listed last above.  This
paper follows the table and counts six, recording the discrepancy rather than
resolving it silently \cite{Yuan2026AEGHoleRegister}.
\end{remark}

\subsection{What the witnesses do not establish}
\label{subsec:p0-holed-nonclaims}

The register's own summary sentence is that a hole "is not a defect to be
removed but the position that carries exchange".  That sentence is the
motivation for this chapter, and it is not proved by any of the six witnesses.
The following list of non-claims is part of the content of this section.

\begin{itemize}
  \item No witness shows that punctures are unavoidable in AES.  The witnesses
  are about chosen finite models: a chosen gluing family, a chosen frame, a
  chosen vocabulary, a chosen stopping rule.  Changing the model changes the
  count, and the register says so for its own semantic witness.
  \item No witness establishes an obstruction of AEG.  The monodromy of the
  frame witness is a monodromy of a stipulated finite flow; the calibration
  witness is a property of a stipulated halting rule; the gluing witness is a
  statement about one gluing family.
  \item The $\mathbb Z_3$ monodromy of the frame witness is not identified with
  any $\mathbb Z_4$ twist of sheet labels appearing elsewhere in the register.
  Those are different groups, and no computation in the register relates them;
  the identification asserted there is recorded as unproved.
  \item The word "hole" in this chapter means a puncture in the sense of
  Definition~\ref{def:p0-punctured-aes}, or an interaction site in the sense of
  the combinatorial witness.  It is not the "one-hole context" of
  Section~\ref{sec:p0-expansion}, which is a placeholder in an expression, and
  the two must not be conflated.
  \item The engineering study that accompanies the register is explicitly
  recorded by its author as an analogy and not as a mathematical assertion; one
  of its three language prototypes was never compiled.  It contributes
  motivation only.
  \item Nothing here changes the status of Paper~I's regular-zero theorems or of
  Paper~III's singular theory.  A punctured AES is a specialization of the
  singular notion, and its classification is not attempted.
\end{itemize}

\subsection{Open problems}
\label{subsec:p0-holed-open}

The register closes with the author's own verdict that this route is not yet
free, for three reviewable reasons: the family of cuts depends on a chosen
permutation, the covering sheets are defined by the cut, and the graded
description depends on an externally supplied frame.  It registers three
directions that have not been attempted, and they are the natural open problems
of this chapter.

\begin{enumerate}
  \item Replace the tangent cone of Proposition~\ref{prop:p0-puncture-tangent-cone}
  by the formal neighbourhood of the puncture: compute the successive jets at
  $p$ and test whether the description ceases to depend on the choice of tangent
  directions.
  \item Blow up the puncture and rewrite the cut-and-glue construction as a
  canonical construction on the exceptional divisor, so that the seam is
  replaced by a construction the data itself determines.
  \item Replace the cut-and-glue calculus by an operadic composition, in which
  the gluing of punctured surfaces is intrinsic and contains no chosen seam.
\end{enumerate}

Each of these is a programme, not a result.  Three further problems arise
directly from the discussion above, and are recorded with their own labels.

\begin{enumerate}
  \item \label{op:p0-four-circle}\emph{Attach AES data to the four-circle
  configuration.}  Construct a surface $\overline M$, a finite set $S$ containing
  the origin, and data $(g,a)$ on $\overline M\setminus S$ such that
  $(\overline M,S,g,a;\mu,\lambda)$ is a punctured AES in the sense of
  Definition~\ref{def:p0-punctured-aes}, such that the four circles (or their
  limiting arcs under $g$) arise naturally from $a$ or from the arithmetic
  motions, and such that the origin is non-extendable.  Until such a
  construction exists, Proposition~\ref{prop:p0-puncture-tangent-cone} is a
  statement about a curve configuration and not about any AES.
  \item \label{op:p0-canonical-punctures}\emph{Make the comparison across a
  puncture canonical.}  Decide whether a relative torsion across a puncture is
  measured by an ACS area, by a holonomy around the puncture, or by neither, and
  in the first case determine the admissible filling classes.
  \item \emph{Make the puncture set canonical.}  Decide whether the set $S$ of
  Definition~\ref{def:p0-punctured-aes} can be determined by the arithmetic data
  rather than stipulated, and whether it is unique when it can.  This is the
  first unfinished item of the exploration register's own programme.
\end{enumerate}

A solution of any of them would turn the language of this chapter from a
description into a theory; the earlier prelude's question, whether a punctured
model carries a ripple structure that survives at the puncture, is a fourth and
independent problem of the same kind.
